\section{Extended Parameters and Benchmarks}
\label{app:ext-params}

This appendix provides extended parameter discussion and additional benchmark data.

\subsection{Admissibility checklist}
A parameter tuple is \emph{admissible} for ARC-SPRUCE if it simultaneously satisfies:
\begin{enumerate}[leftmargin=1.25em]
  \item \textbf{Bound vs.\ modulus:} \(\BoundMsg \ll q\) so membership checks on
  \([-\BoundMsg,\BoundMsg]\) do not wrap modulo \(q\).
  \item \textbf{Commitment hiding:} the commitment-only randomness \(\bar r\) makes the factorization-aware
  bound of Lemma~\ref{lem:ajtai-hiding} at most \(2^{-\lambda}\).
  \item \textbf{Target hiding:} the long BB-tran randomizer \(r_0\in\Rq^{\ell_{r_0}}\) makes the
  factorization-aware bound of Lemma~\ref{lem:htran-lhl} at most \(2^{-\lambda}\).
  \item \textbf{Sampler interface:} the trapdoor/preimage layer satisfies Definition~\ref{def:trapdoor-suite},
  including the public \(\ell_\infty\) shortness bound \(\BoundSig\).
  \item \textbf{Rational-hash guard:} denominator invertibility failure is negligible or explicitly carried as the
  rational-hash admissibility error \(\delta\).
  \item \textbf{SmallWood degree/soundness:} the DECS/LVCS/PCS/PIOP degree bounds and soundness parameters are
  computed for the actual issuance and showing relations.
  \item \textbf{PRF key space and tag length:} the relation enforces
  \(\veck\in\mathcal K_{\mathsf{key}}\), the PRF assumption is made for this same key space, and
  \(\ell_{\prftag}\) is large enough for the intended tag length.
\end{enumerate}

\subsection{Smoothing and sampler budgets}
Let \(\eta_\varepsilon(\Lambda)\) denote the smoothing parameter of a lattice \(\Lambda\). A sufficient condition for
statistical hiding in Gaussian preimage sampling is to choose Gaussian widths above a floor
\(\sigma_{\min}\ge \eta_\varepsilon(\Lambda_\perp(A))\), with slack (Appendix~\ref{app:gaussian-sampler}).
For the ambient lattice \(\Z^{2N}\) one may use the classical bound
\[
  \eta_\varepsilon(\Z^{2N})
  \,=\,
  \sqrt{\ln\!\bigl(4N(1+2^\lambda)\bigr)/\pi}.
\]
In our tuning, trapdoor quality \(\alpha\) and acceptance parameters are chosen so that the output shortness bound
\(\BoundSig\) holds with overwhelming probability.

\subsection{Example lattice parameters}
Table~\ref{tab:lattice-params} records one concrete instantiation compatible with the sizing discussion in
Section~\ref{sec:parameters}.

\begin{table}[H]
\centering
\renewcommand{\arraystretch}{1.15}
\begin{tabular}{@{}lll@{}}
\toprule
\textbf{Knob / Result} & \textbf{Symbol} & \textbf{Value (example)} \\
\midrule
Ring degree & \(N\) & \(1024\) \\
Modulus & \(q\) & \(1038337\) \\
Security target & \(\lambda\) & \(128\) \\
\midrule
Trapdoor quality & \(\alpha\) & \(\approx 1.23\) (measured) \\
Gaussian width scale & \(\sigma\) & tuned per-slot (Appendix~\ref{app:gaussian-sampler}) \\
Signature bound & \(\BoundSig\) & derived from sampler tail bound \\
\bottomrule
\end{tabular}
\caption{Example lattice parameters (illustrative; see Appendix~\ref{app:gaussian-sampler} for the sampling interface and
tuning).}
\label{tab:lattice-params}
\end{table}

\subsection{LHL design points for the long $r_0$-randomizer}
For the measured parameter scale \(N=1024\), \(q=1038337\), and \(\lambda=128\), note that
\[
  q\equiv 1 \pmod {2N}.
\]
Thus \(X^N+1\) splits completely over \(\F_q\), so \(d_{\mathrm{fac}}=1\) and
\(f_{\mathrm{fac}}=N\).  For a one-coordinate target (\(n=1\)), Lemma~\ref{lem:htran-lhl} gives
\[
  \varepsilon_{r_0}(\beta_{r_0},\ell_{r_0})
  =
  \frac12
  \sqrt{
    \left(
      1+\frac{q}{(2\beta_{r_0}+1)^{\ell_{r_0}}}
    \right)^N
    -1
  }.
\]
The target-hiding condition is
\[
  \varepsilon_{r_0}(\beta_{r_0},\ell_{r_0})\le 2^{-\lambda}.
\]
Equivalently, in the small-error approximation one needs roughly
\[
  \ell_{r_0}\log_2(2\beta_{r_0}+1)
  \gtrsim
  \log_2 q+2\lambda+\log_2 N-2,
\]
not \(\log_2 q+2\lambda/N\).

Table~\ref{tab:lhl-design-points} records representative design points computed from the displayed
factorization-aware bound.

\begin{table}[H]
\centering
\renewcommand{\arraystretch}{1.15}
\begin{tabular}{@{}ccc@{}}
\toprule
\(\beta_{r_0}\) & minimum \(\ell_{r_0}\) for \(\lambda=128\) & membership degree \\
\midrule
1  & 180 & 3  \\
5  & 83  & 11 \\
8  & 70  & 17 \\
31 & 48  & 63 \\
\bottomrule
\end{tabular}
\caption{Representative target-hiding design points for the fully split example
\(N=1024,q=1038337,n=1\).  These values use the factorization-aware cyclotomic LHL bound.}
\label{tab:lhl-design-points}
\end{table}

The table shows that the fully split example requires a much wider \(r_0\) than a total-entropy-only
calculation would suggest.  Increasing \(\beta_{r_0}\) reduces the required width \(\ell_{r_0}\), but it increases
the degree of the range-check polynomial.

\subsection{SmallWood soundness terms}
SmallWood’s soundness is a combination of DECS polynomial binding, LVCS/PCS binding, PIOP sampling errors,
and Fiat--Shamir/grinding terms~\cite{EPRINT:FenRiv25b}.  For parameters
\((N,s,\ell,\ell',\rho,\eta)\), field \(\F\), and degree bounds \(d_{\mathrm{decs}}=s+\ell-1\) and \(d_Q\),
schematic source-compatible terms include
\[
\varepsilon_{\mathrm{decs}}
=
\binom{N}{d_{\mathrm{decs}}+2}|\F|^{-\eta}
+
\frac{Q_{\RO}^2}{2^{2\lambda}},
\qquad
\varepsilon_{\mathrm{sum}} = |\F|^{-\rho},
\qquad
\varepsilon_{\mathrm{tail}} =
\frac{\binom{d_Q}{\ell'}}{\binom{|S|}{\ell'}}.
\]
The final knowledge-soundness bound for ARC-SPRUCE must instantiate the exact SmallWood-ARK theorem with
the actual row inventory, the actual \(d_Q\), the Fiat--Shamir query bounds, and the grinding parameters.
Our issuance statement includes the bilinear inverse check \((B_3-\mathbf 1_n r_1)\Had Z=1\) together with
the linear target identity given public \(T\), so \(d_Q\) must account for the inverse relation and the
range/membership checks.  The showing statement additionally depends on the PRF effective degree
(Appendix~\ref{app:prf}).

\subsection{SmallWood knobs (example)}
For the issuance and showing proofs discussed in Section~\ref{sec:parameters}, one example parameter tuple is:
\[
  (s,\ell,\ell',\rho,\theta,\eta)\;=\;(16,25,2,2,4,19),
\]
with masking degree \(d_Q\) derived from the instantiated constraint degrees. The pre-sign statement includes the bilinear inverse-witness constraint together with linear commitment, centering, and target equations, while the post-sign statement adds the PRF S-box composition on top of the same BB-tran inverse check.

\subsection{Extended benchmark grid (SmallWood)}
For additional context on how SmallWood knobs affect proof size and proving time, Table~\ref{tab:results-grid}
reports an example sweep over packing and soundness parameters for standalone SmallWood-style statements.  These
figures are not a final ARC-SPRUCE proof-size claim unless the row inventory, constraint degrees, PRF checkpoint
schedule, and Fiat--Shamir soundness terms are instantiated for the exact ARC relations.

\begin{table}[H]
\centering
\footnotesize
\setlength{\tabcolsep}{4pt}
\begin{tabular}{r r r r r r r r r}
\toprule
\textbf{ProofKB} & \textbf{Time (s)} & \textbf{Bits} & \textbf{NCols} & \(\ell\) & \(\ell'\) & \(\rho\) & \(\theta\) & \(\eta\)\\
\midrule
\ \ 9.32 & 0.266 & 133.44 & 6 & 22 & 2 & 2 & 4 & 17 \\
10.07 & 0.228 & 146.93 & 4 & 24 & 2 & 2 & 4 & 17 \\
11.52 & 0.250 & 133.22 & 4 & 28 & 2 & 2 & 4 & 17 \\
\midrule[1.2pt]
14.06 & 0.630 & 198.01 & 4 & 34 & 8 & 5 & 2 & 26 \\
14.17 & 0.654 & 192.26 & 6 & 34 & 8 & 5 & 2 & 26 \\
14.24 & 0.670 & 198.48 & 4 & 34 & 8 & 6 & 2 & 26 \\
\bottomrule
\end{tabular}
\caption{Example sweep over SmallWood knobs for standalone parameter exploration.  Sizes exclude public
parameters and are not a final ARC-SPRUCE proof-size claim.}
\label{tab:results-grid}
\end{table}
